\documentclass[11pt,letterpaper]{article}
\usepackage[T1]{fontenc}
\usepackage{lmodern}
\usepackage[margin=1in]{geometry}
\usepackage{microtype}
\usepackage{amsmath,amssymb,amsthm,mathtools}
\usepackage{enumitem}
\usepackage{xcolor}
\usepackage{hyperref}
\hypersetup{colorlinks=true,linkcolor=blue!45!black,citecolor=blue!45!black,
  urlcolor=blue!45!black,pdftitle={Conway subprime closure: weaker analytic inputs},
  pdfauthor={OpenAI Astra; Romain Popescu},
  pdfsubject={A conditional reduction using prime-scale exceptional sets and proportional prime intervals}}
\setlist[enumerate]{leftmargin=*,itemsep=0.65em,topsep=0.6em}
\numberwithin{equation}{section}
\newtheorem{theorem}{Theorem}[section]
\newtheorem{lemma}[theorem]{Lemma}
\newtheorem{proposition}[theorem]{Proposition}
\theoremstyle{remark}
\newtheorem{remark}[theorem]{Remark}
\DeclareMathOperator{\lpf}{lpf}
\DeclareMathOperator{\meas}{meas}
\newcommand{\N}{\mathbb N}
\newcommand{\Z}{\mathbb Z}
\newcommand{\R}{\mathbb R}
\newcommand{\Pp}{\mathbb P}
\newcommand{\ph}{\varphi}
\newcommand{\tot}{\phi_{\!E}}
\newcommand{\ind}{\mathbf 1}
\newcommand{\ee}[1]{e\!\left(#1\right)}
\newcommand{\floor}[1]{\lfloor #1\rfloor}
\newcommand{\ceil}[1]{\lceil #1\rceil}
\newcommand{\HH}[2]{H_{#1}(#2)}
\allowdisplaybreaks[2]
\title{Conway's subprime closure:\\a reduction with weaker analytic inputs}
\author{OpenAI Astra \and Romain Popescu}
\date{8 September 2026\\[3pt]\small Research variant: weaker inputs}


\begin{document}
\maketitle
\vspace{-1.2em}

\section{Problem and scope of the revision}

Let $s(1)=1$, let $s(m)=m$ for prime $m$, and let
$s(m)=m/\lpf(m)$ for composite $m$. Define
\[
 C_0=\{1\},\qquad C_{n+1}=C_n\cup\{s(a+b):a,b\in C_n\},
 \qquad M_n=\max C_n.
\]
Repeated inputs are allowed. Caragiu, Vicol, and Zaki
\cite[Theorem~1 and Conjecture~3]{CVZ} prove that $\bigcup_n C_n=\N$
and ask whether $|C_{n+1}|/|C_n|\to\ph=(1+\sqrt5)/2$.

This is a separate revision of the preserved proof. Its purpose is to
reduce the analytic statements needed for a Lean formalization. We prove
an explicit \emph{conditional reduction}: ordinary proportional-interval
prime estimates and restricted binary representation estimates with
$o(X/\log X)$ exceptions suffice. The full logarithmic savings
$O_A(X/\log^A X)$, prime estimates in shrinking intervals, and the
almost-equal-summands theorem are not hypotheses of this reduction.
The analytic hypotheses below are stated, not formalized or discharged here.

\begin{theorem}\label{thm:weak-main}
Assume the prime estimates and the two restricted binary hypotheses of
Section~\ref{sec:inputs}. Then
\begin{equation}\label{eq:weak-limits}
 \frac{M_{n+1}}{M_n}\longrightarrow\ph,\qquad
 \frac{|C_n|}{M_{n-1}}\longrightarrow1,\qquad
 \frac{|C_{n+1}|}{|C_n|}\longrightarrow\ph,\qquad
 \frac{|C_n|}{M_n}\longrightarrow\frac1\ph.
\end{equation}
Moreover, for each $0<\zeta<1$, every prime at most $(1-\zeta)M_n$
belongs to $C_n$ for all sufficiently large $n$. Consequently the complete
prime prefix $Q_n$ of the preserved proof satisfies $Q_n/M_n\to1$.
\end{theorem}

The revised argument does not assert a positive limit for $M_n/\ph^n$.
That stronger conclusion is unnecessary for \eqref{eq:weak-limits}.
The proof is a research draft, not an independently reviewed or fully
Lean-checked theorem.

\section{Proof overview}

\begin{enumerate}
\item \textbf{Keep the elementary upper bound.} Composite sums give outputs
at most the preceding maximum. Parity therefore gives
$M_{n+1}\le M_n+M_{n-1}$. The weighted quantity
$W_j=M_{j+1}+M_j/\ph$ grows by at most a factor $\ph$ per generation.

\item \textbf{Propagate a slightly slower prime profile.} Fix
$\lambda<\ph$, as close to $\ph$ as needed. If complete prime cutoffs
$L,\lambda L$ are present in consecutive generations, buffered Goldbach
filling and a prime-counting pigeonhole propagate the cutoffs to
$\lambda^2L,\lambda^3L,\ldots$. All intervals have fixed proportional
lengths; no errors need to be summed over generations.

\item \textbf{Increase the profile when a maximum gets ahead.} An available
prime $t>(1+\delta)L$ gives representations
$2p+q=2u-t$, with $p,q$ in known prime intervals. Almost every candidate
prime $u$ is then generated. Pairing these $u$ with integers in a buffered
filled interval extends a later complete prime cutoff by a fixed amount.
After three generations the entire two-cutoff profile can be increased
by a fixed factor $1+\kappa>1$.

\item \textbf{Use a bounded-step comparison.} Compare $W_j$ with the
weighted profile $(\lambda+\ph^{-1})L$. A normal step increases their
ratio by at most $\ph/\lambda$. If that ratio is too large, the increase
of the profile gives a strict contraction within three or four steps.
Taking $\lambda$ sufficiently close to $\ph$ forces the profile to stay
close to the maxima, and hence proves approximate prime completeness.

\item \textbf{Take the ratio limit directly.} Prime completeness makes
the Fibonacci upper bound asymptotically an equality. The resulting
ratio recurrence is a contraction toward $\ph$. Buffered filling then
gives $|C_n|\sim M_{n-1}$ and the desired cardinality limit.
\end{enumerate}

\section{Precisely weakened analytic inputs}\label{sec:inputs}

Let $\Pp$ be the set of primes. All logarithms are natural; all integer
counts in real intervals use their intersection with $\N$. Put
\[
 \mathcal L(X)=X/\log X\quad(X>1),\qquad
 \operatorname{Pref}(j,L)\ :\Longleftrightarrow\
       \forall p\in\Pp,\ p\le L\Longrightarrow p\in C_j,
\]
\[
 H_j(Y)=\#\bigl(\{1,\ldots,\floor Y\}\setminus C_j\bigr).
\]
The notation $f(X)=o(\mathcal L(X))$ means
$f(X)/\mathcal L(X)\to0$ as $X\to\infty$.

\paragraph{Prime estimates.}
We use only the following consequences of the ordinary prime number theorem:
for each fixed $0<\alpha<\beta$ there is $c_{\alpha,\beta}>0$ such that
\begin{equation}\label{eq:prime-lower}
 \#(\Pp\cap[\alpha X,\beta X])
 \ge c_{\alpha,\beta}\mathcal L(X)
 \quad\text{for all sufficiently large }X,
\end{equation}
and
\begin{equation}\label{eq:prime-density}
 \pi(X)=o(X).
\end{equation}
In particular, for each fixed $0<\rho<1$ and sufficiently large $L$,
\begin{equation}\label{eq:prefix-max}
 \operatorname{Pref}(j,L)\Longrightarrow
          M_j\ge(1-\rho)L.
\end{equation}
Indeed \eqref{eq:prime-lower} supplies a prime in $[(1-\rho)L,L]$.
The threshold is independent of $j$. These estimates do not involve a
prime-number-theorem error term.

\paragraph{Restricted binary hypotheses.}
For $\nu\in\{1,2\}$ and fixed $0<\alpha<\beta$, $0<\gamma<\delta_1$,
$\eta>0$, $B>0$, define
\[
 I_X=[\alpha X,\beta X],\quad J_X=[\gamma X,\delta_1X],\quad
 \mathcal J_{\nu,X}(N)=\meas\{v\in I_X:N-\nu v\in J_X\}.
\]
Let $E_\nu(X)$ count the integers $N$ satisfying
\[
 1\le N\le BX,\qquad N\equiv\nu+1\pmod2,\qquad
 \mathcal J_{\nu,X}(N)\ge\eta X,
\]
which have no representation $N=\nu p+q$ with primes $p\in I_X$, $q\in J_X$.
The hypothesis, separately for $\nu=1$ and $\nu=2$, is
\begin{equation}\label{eq:binary-weak}
 E_\nu(X)=o(\mathcal L(X)).
\end{equation}
Every interval parameter is fixed before the limit is taken. We assume no
uniform rate as those parameters vary. In the proof only finitely many
intervals are needed after fixing the desired accuracy.

\begin{remark}\label{rem:rate-distinction}
The rate in \eqref{eq:binary-weak} is strictly stronger than density zero,
which says $E_\nu(X)=o(X)$. For example, the primes have density zero but
number $\asymp X/\log X$. An exceptional set of that size can contain
\emph{every} prime candidate in one of our pigeonholes. Nor does an
unrestricted Goldbach representation theorem by itself provide a
representation with summands in specified intervals. Thus a theorem stated
only as \texttt{DensityZero} for unrestricted Goldbach exceptions is not
a direct substitute for \eqref{eq:binary-weak}.
\end{remark}

\section{Buffered filling and propagation}

\begin{lemma}[Buffered filling]\label{lem:buffer}
Assume \eqref{eq:binary-weak} for $\nu=1$. For every fixed $0<\theta<1$
there is a function $e_\theta(X)\to0$ such that, for all sufficiently large
$X$ and every $j$,
\begin{equation}\label{eq:buffer}
 \operatorname{Pref}(j,X)\Longrightarrow
 H_{j+1}((1-\theta)X)\le e_\theta(X)\mathcal L(X).
\end{equation}
\end{lemma}

\begin{proof}
For a scale $Y$, take both prime intervals to be
$[(1-\theta)Y/4,Y]$. If $(1-\theta)Y/2\le m\le(1-\theta)Y$,
the real overlap for $2m=p+q$ is at least
\[
 \eta_\theta Y,\qquad
 \eta_\theta=\min\bigl((1-\theta)/2,2\theta\bigr)>0.
\]
To check this, write $a=(1-\theta)Y/4$. The overlap of
$[a,Y]$ and $[2m-Y,2m-a]$ has length
$\min(2m-2a,2Y-2m,Y-a)$, and each of these three terms is at least
$\eta_\theta Y$ in the stated range.
Thus \eqref{eq:binary-weak} with $B=2$ gives at most
$e(Y)Y/\log Y$ exceptional midpoints, for a function $e(Y)\to0$.
Every other midpoint $m\ge2$ is generated by $s(p+q)=s(2m)=m$
once every prime up to $Y$ is present.

Apply this at the scales $Y_r=X/2^r$ with $Y_r\ge\sqrt X$.
Their midpoint ranges cover all of
$[\sqrt X,(1-\theta)X]$ except possibly a bottom interval of length
$O(\sqrt X)$. More explicitly the union starts below $\sqrt X$,
since the last scale is less than $2\sqrt X$.
Also $\sum_rY_r\le2X$ and $\log Y_r\ge\tfrac12\log X$. Hence
\[
 \sum_r e(Y_r)\frac{Y_r}{\log Y_r}
 \le 4\sup_{Y\ge\sqrt X}e(Y)\frac{X}{\log X}
 =o(\mathcal L(X)).
\]
The omitted $O(\sqrt X)$ integers are also $o(\mathcal L(X))$.
Under $\operatorname{Pref}(j,X)$ every scale uses available primes.
The bounds depend on $X,\theta$, not on $j$, proving the uniform assertion.
\end{proof}

\begin{lemma}[Proportional extension]\label{lem:prop-extension}
Fix $0<\epsilon<1$ and $B\ge1$. For sufficiently large $Y$, if
$Y\le X\le BY$, $\operatorname{Pref}(j,Y)$ and
$\operatorname{Pref}(j+1,X)$, then
\begin{equation}\label{eq:prop-extension}
 \operatorname{Pref}\bigl(j+2,(1-\epsilon)(X+Y)\bigr).
\end{equation}
\end{lemma}

\begin{proof}
Only a prime $r$ with $X<r\le(1-\epsilon)(X+Y)$ needs attention;
if this range is empty the assertion follows by retention. For each prime
$p\in[(1-\epsilon)X,X]$, the positive integer $a=r-p$ is at most
$(1-\epsilon)Y$. The candidates $a$ are distinct. There are
$\gg_\epsilon X/\log X$ candidates by \eqref{eq:prime-lower}, but
only $o(Y/\log Y)=o(X/\log X)$ missing integers in that buffered range
of $C_{j+1}$, by Lemma~\ref{lem:buffer} and $Y\le X\le BY$.
Some $a$ is present; $p$ is present by the second prefix hypothesis.
Since $r$ is prime, $s(p+a)=r\in C_{j+2}$.
\end{proof}

\section{A finite increase of a prime profile}\label{sec:boost}

Fix for this section
\begin{equation}\label{eq:parameters}
 0<\delta\le\frac1{10},\qquad \frac85\le\lambda<\ph,\qquad
 \epsilon=1-\frac{\lambda^2}{\lambda+1},\qquad
 a=\frac{\delta}{100},\quad b=\frac{\delta}{8},\quad
 \kappa=\frac{(1-\epsilon)b}{\lambda^4}.
\end{equation}
Then $0<\epsilon\le1/65$ and
\begin{equation}\label{eq:lambda-recursion}
 \lambda^2=(1-\epsilon)(\lambda+1).
\end{equation}
A \emph{profile at $(j,L)$} means
$\operatorname{Pref}(j,L)$ and $\operatorname{Pref}(j+1,\lambda L)$.
Lemma~\ref{lem:prop-extension} implies that a profile propagates to
$(j+1,\lambda L)$, and hence to $(j+r,\lambda^rL)$ for every $r\ge0$,
once $L$ exceeds a fixed threshold. All ratios of adjacent cutoffs here
are $\lambda<2$.

\begin{lemma}[Increase of the profile]\label{lem:boost}
For the fixed parameters \eqref{eq:parameters} there is a threshold $L_0$
such that, uniformly in $j$, a profile at $(j,L)$ with $L\ge L_0$ and
\begin{equation}\label{eq:outlier-condition}
 M_j>(1+\delta)L
\end{equation}
produces a profile at
\begin{equation}\label{eq:boosted-pair}
 \bigl(j+3,(1+\kappa)\lambda^3L\bigr).
\end{equation}
\end{lemma}

\begin{proof}
The elementary facts used here are proved in Lemma~\ref{lem:elementary}
below. Choose the first index $i$ for which $M_i>(1+\delta)L$.
For large $L$, $i\ge1$, and \eqref{eq:outlier-condition} implies $i\le j$.
The bound $M_i\le2M_{i-1}$ gives an available odd prime
\begin{equation}\label{eq:t-range}
 t=M_i\in C_j,\qquad 1+\delta<t/L\le2+2\delta.
\end{equation}
It is odd once $L$ is large because $t>2$.

Choose a finite set of real numbers $\mathcal T$ such that every
$\tau\in[1+\delta,2+2\delta]$ is within $a/4$ of some $T\in\mathcal T$,
and all its members lie in
$[1+\delta-a/4,2+2\delta+a/4]$. A grid of mesh at most $a/2$
suffices. Select $T$ with $|t/L-T|\le a/4$, and put
\begin{equation}\label{eq:boost-intervals}
 \begin{aligned}
 I_L&=[(\lambda-2a)L,(\lambda-a)L],\\
 J_L&=[(1-5a)L,(1-a/2)L],\\
 U_{T,L}&=[(\lambda+(T+1)/2-3a)L,
            (\lambda+(T+1)/2-5a/2)L].
 \end{aligned}
\end{equation}
These are fixed positive intervals after $\lambda,\delta,T$ are chosen.
Every prime in $I_L$ belongs to $C_{j+1}$ and every prime in $J_L$
belongs to $C_j$.

For an integer $u\in U_{T,L}$, the target $N=2u-t$ is odd. For every
real $p\in I_L$, the value $q=2u-t-2p$ satisfies
\begin{equation}\label{eq:q-boost-range}
 (1-17a/4)L\le q\le(1-3a/4)L.
\end{equation}
Indeed the endpoint bounds before using $|T-t/L|\le a/4$ are
$(1-4a+T-t/L)L$ and $(1-a+T-t/L)L$.
The interval in \eqref{eq:q-boost-range} lies inside $J_L$, so the
real overlap for $2p+q=N$ has length $aL$.
Also $0<N\le6L$ for large $L$: the endpoint coefficients for $N$
lie between $2\lambda+1-25a/4$ and $2\lambda+1-19a/4$,
and $8/5\le\lambda<5/3$, $a\le1/1000$.

Apply \eqref{eq:binary-weak} with $\nu=2$, $\eta=a$, $B=6$.
There are only finitely many $T$, so the maximum of their error functions
still tends to zero. The injective map $u\mapsto2u-t$ shows that all but
$o(L/\log L)$ integers $u\in U_{T,L}$ admit the required prime representation,
uniformly over $j,t,T$. If such a $u$ is itself prime, then
\begin{equation}\label{eq:boost-operations}
 \frac{t+q}{2}=s(t+q)\in C_{j+1},\qquad
 u=p+\frac{t+q}{2}=s\left(p+\frac{t+q}{2}\right)\in C_{j+2}.
\end{equation}
Here $t,q$ are odd primes, so their sum is composite with least prime
factor $2$. We have proved
\begin{equation}\label{eq:u-exceptions}
 \#\bigl((\Pp\cap U_{T,L})\setminus C_{j+2}\bigr)
       =o(L/\log L).
\end{equation}

We next show that every prime up to $(\lambda^3+b)L$ belongs to $C_{j+3}$.
Ordinary propagation already supplies every prime at most $\lambda^3L$.
Let $r$ be a prime in the remaining range
$[\lambda^3L,(\lambda^3+b)L]$, and let $u\in\Pp\cap U_{T,L}$.
For $v=r-u$, the endpoint calculations give
\begin{equation}\label{eq:v-boost-range}
 \frac12 L\le v\le(\lambda-11\delta/32)L
                     <(\lambda-\delta/4)L.
\end{equation}
For the upper bound, use $T\ge1+\delta-a/4$ and
$\lambda^3-2\lambda-1\le0$ for $8/5\le\lambda\le\ph$:
\begin{align*}
 v/L
 &\le\lambda^3+b-\lambda-(T+1)/2+3a\\
 &\le\lambda+b-\delta/2+25a/8
 =\lambda-11\delta/32.
\end{align*}
For the lower bound, $\lambda^3-\lambda\ge312/125$ and
$T\le2+2\delta+a/4$ yield
\[
 v/L\ge\lambda^3-\lambda-(T+1)/2+5a/2
 \ge\frac{249}{250}-\delta+\frac{19a}{8}
 \ge\frac{112}{125}>\frac12.
\]
The two polynomial bounds follow because $x^3-2x-1$ and $x^3-x$
are increasing for $x\ge8/5$, and $\ph^3-2\ph-1=0$.

Use Lemma~\ref{lem:buffer} with prime cutoff $\lambda L$ in $C_{j+1}$
and buffer $\theta=\delta/(4\lambda)$. It bounds the missing integers
of $C_{j+2}$ up to $(\lambda-\delta/4)L$ by $o(L/\log L)$.
There are $\gg_{\lambda,\delta,T}L/\log L$ prime candidates in $U_{T,L}$,
whose length is $aL/2$. Exclude those missing by \eqref{eq:u-exceptions}
and those with $r-u\notin C_{j+2}$. The latter exclusion is bounded by
the buffered hole count because $u\mapsto r-u$ is injective.
A candidate remains. Its two inputs generate $r$ in $C_{j+3}$.
The thresholds are uniform in $r$, so this proves the entire asserted prefix.

At time $j+2$ the ordinary prefix is $\lambda^2L$; at time $j+3$ the
improved prefix is $(\lambda^3+b)L$. Their ratio is less than $2$.
One more application of Lemma~\ref{lem:prop-extension} gives prefix
\[
 (1-\epsilon)(\lambda^3+b+\lambda^2)L
   =\bigl(\lambda^4+(1-\epsilon)b\bigr)L
   =(1+\kappa)\lambda^4L
\]
at time $j+4$. At time $j+3$ the improved cutoff also exceeds
$(1+\kappa)\lambda^3L$, since
$\kappa\lambda^3=(1-\epsilon)b/\lambda\le b$.
These are exactly the two cutoffs in \eqref{eq:boosted-pair}.
\end{proof}

\section{The bounded-step comparison}

\begin{lemma}[Elementary upper bounds]\label{lem:elementary}
For $n\ge1$, $M_n$ is prime and
\begin{equation}\label{eq:elementary}
 C_n\subseteq[1,M_{n-1}]\cup\Pp,\qquad
 M_{n+1}\le M_n+M_{n-1},\qquad M_n\le2M_{n-1}.
\end{equation}
With $W_j=M_{j+1}+M_j/\ph$ we have $W_{j+1}\le\ph W_j$.
\end{lemma}

\begin{proof}
A composite sum of two inputs at most $M_{n-1}$ gives an output at most
$M_{n-1}$. Any new element above that maximum is therefore prime.
Since $M_1=2$ and $M_2=3$, all maxima from $M_1$ on are prime, and all
from $M_2$ on are odd. Every even element of $C_n$ for $n\ge2$ is at
most $M_{n-1}$, including the even prime $2$. A new maximum at time
$n+1$ is an odd prime formed from an odd and an even input, so it is
at most $M_n+M_{n-1}$. The case $n=1$ is $3\le2+1$; a retained
maximum also satisfies the inequality. Monotonicity gives
$M_n\le2M_{n-1}$, with the first case checked directly. Finally
\[
 W_{j+1}\le M_{j+1}+M_j+M_{j+1}/\ph
 =\ph M_{j+1}+M_j=\ph W_j.
\]
\end{proof}

\begin{proposition}[Approximate prime completeness]\label{prop:prefix-close}
For each $0<\zeta<1$, $\operatorname{Pref}(n,(1-\zeta)M_n)$ holds for
all sufficiently large $n$.
\end{proposition}

\begin{proof}
Fix $0<\delta\le1/10$ and make the explicit choice
\[
 \lambda=\ph-\frac{\delta}{100}\left(\ph-\frac85\right).
\]
Then $8/5\le\lambda<\ph$. With the parameters of \eqref{eq:parameters},
this choice satisfies
\begin{equation}\label{eq:choose-lambda}
 A_0:=\frac\ph\lambda>1,\qquad A_0^5\le1+\delta,\qquad
 \rho:=\frac{A_0^4}{1+\kappa}<1.
\end{equation}
Here is a direct verification. Since $8/5<\ph<5/3$, the number
$y=A_0-1$ satisfies $0<y\le\delta/2400<1/100$.
Expanding powers gives $(1+y)^4\le1+5y$ and
$(1+y)^5\le1+6y$ for $0\le y\le1/100$.
Also $1-\epsilon\ge1/2$ and $\lambda^4<8$, so
$\kappa\ge\delta/128$. Therefore
$A_0^4\le1+\delta/480<1+\kappa$ and
$A_0^5\le1+\delta/400\le1+\delta$, as required.
Only this fixed $\lambda$ is used in the rest of the argument.

Take $L$ beyond all thresholds in the propagation and increase lemmas,
and beyond the threshold for \eqref{eq:prefix-max} with accuracy $\delta$.
Exhaustion gives a $j$ with $[1,\ceil{\lambda L}]\subseteq C_j$,
so there is a profile at $(j,L)$. For any profile, define
\[
 Z(j,L)=\frac{M_{j+1}+M_j/\ph}{(\lambda+1/\ph)L}.
\]
A normal one-step propagation satisfies
\begin{equation}\label{eq:normal-Z}
 Z(j+1,\lambda L)\le A_0 Z(j,L).
\end{equation}
If $Z(j,L)>1+\delta$, at least one of
\[
 M_j>(1+\delta)L,\qquad M_{j+1}>(1+\delta)\lambda L
\]
holds; otherwise their weighted sum contradicts the strict inequality
for $Z$. In the first case Lemma~\ref{lem:boost} gives an increased
profile after $h=3$ steps. In the second case perform one normal step
and then use that lemma, giving an increased profile after $h=4$ steps.
In either case the new pair is
$(j+h,(1+\kappa)\lambda^hL)$ and
\begin{equation}\label{eq:boost-Z}
 Z\bigl(j+h,(1+\kappa)\lambda^hL\bigr)
 \le\frac{A_0^h}{1+\kappa}Z(j,L)\le\rho Z(j,L).
\end{equation}

Construct a sequence of profiles as follows. If $Z\le1+\delta$, take
one normal step; if $Z>1+\delta$, take the increase of length $3$ or $4$
just described. All scales increase, so all thresholds remain valid.
The generation indices increase by integers between $1$ and $4$ and
hence tend to infinity. There is a profile in this sequence with
$Z\le1+\delta$: otherwise \eqref{eq:boost-Z} would give
$Z_m\le\rho^mZ_0\to0$, contradicting $Z_m>1+\delta$.
After the first such profile, every subsequent selected profile satisfies
\begin{equation}\label{eq:node-bound}
 Z\le A_0(1+\delta).
\end{equation}
This is an induction: a normal step starts with $Z\le1+\delta$ and
uses \eqref{eq:normal-Z}; an increase contracts the preceding bound.

We must also control generations between selected profiles. From a
selected pair $(j,L)$, normal propagation supplies a profile
$(j+r,\lambda^rL)$ at every intermediate index. For $0\le r<h\le4$,
\[
 Z(j+r,\lambda^rL)\le A_0^r Z(j,L)
 \le A_0^5(1+\delta)\le(1+\delta)^2.
\]
(The exponent $5$ is a harmless upper bound.) Therefore, at every
sufficiently large index $n$, there is an $L_n$ such that
\begin{equation}\label{eq:all-stage-profile}
 \operatorname{Pref}(n,L_n),\qquad
 \operatorname{Pref}(n+1,\lambda L_n),\qquad
 Z(n,L_n)\le(1+\delta)^2.
\end{equation}
These scales tend to infinity: selected scales grow by at least $\lambda$
per elapsed generation and intermediate scales are their corresponding
normal propagations.

By \eqref{eq:prefix-max}, $M_{n+1}\ge(1-\delta)\lambda L_n$.
Solving the last inequality of \eqref{eq:all-stage-profile} for $M_n$ gives
\begin{align*}
 \frac{M_n}{L_n}
 &\le(1+\delta)^2(\ph\lambda+1)-(1-\delta)\ph\lambda\\
 &=1+(2\delta+\delta^2)(\ph\lambda+1)+\delta\ph\lambda
 \le1+12\delta.
\end{align*}
For the last inequality, use $\delta\le1/10$,
$\ph\lambda<3$, and $\ph\lambda+1<4$; the excess is at most
$(2.1\cdot4+3)\delta=11.4\delta$.
It follows that $\operatorname{Pref}(n,M_n/(1+12\delta))$ holds eventually.
Given $0<\zeta<1$, choose $0<\delta\le\min(1/10,\zeta/24)$.
Since $1/(1+12\delta)\ge1-12\delta\ge1-\zeta$, this proves the proposition.
\end{proof}

\section{Ratios and cardinalities without positive amplitudes}

\begin{proof}[Proof of Theorem~\ref{thm:weak-main}]
Exhaustion implies $M_n\to\infty$. Fix $0<\zeta<1$. Proposition~\ref{prop:prefix-close}
gives consecutive prime cutoffs
$X=(1-\zeta)M_n$ and $Y=(1-\zeta)M_{n-1}$ for all large $n$.
They satisfy $Y\le X\le2Y$. Lemma~\ref{lem:prop-extension}, with
$\epsilon=\zeta$, gives every prime up to
$(1-\zeta)^2(M_n+M_{n-1})$ in $C_{n+1}$.
Using \eqref{eq:prefix-max} once more,
\[
 (1-\zeta)^3(M_n+M_{n-1})\le M_{n+1}\le M_n+M_{n-1}
\]
for all sufficiently large $n$. Since $\zeta$ is arbitrary,
\begin{equation}\label{eq:fib-saturation}
 \frac{M_{n+1}}{M_n+M_{n-1}}\longrightarrow1.
\end{equation}

Put $r_n=M_{n+1}/M_n$. By \eqref{eq:elementary}, $1\le r_n\le2$.
Equation~\eqref{eq:fib-saturation} implies
$r_n=1+1/r_{n-1}+e_n$, where $e_n\to0$: explicitly $e_n$ is
$1+1/r_{n-1}$ times the difference between the quotient in
\eqref{eq:fib-saturation} and $1$, and that multiplier is at most $2$.
Since $\ph=1+1/\ph$,
\[
 |r_n-\ph|\le\frac{|r_{n-1}-\ph|}{\ph r_{n-1}}+|e_n|
 \le\ph^{-1}|r_{n-1}-\ph|+|e_n|.
\]
The finite number $D=\limsup_n|r_n-\ph|$ satisfies $D\le\ph^{-1}D$,
and hence $D=0$. Thus $r_n\to\ph$.

For the cardinality, again fix $0<\zeta<1$. Prime completeness and
Lemma~\ref{lem:buffer} show
\[
 H_n((1-\zeta)^2M_{n-1})=o(M_{n-1}/\log M_{n-1}).
\]
Consequently
\[
 \liminf_n\frac{|C_n|}{M_{n-1}}\ge(1-\zeta)^2.
\]
The upper structural bound and \eqref{eq:prime-density} give
\[
 |C_n|\le M_{n-1}+\pi(M_n)
       \le M_{n-1}+\pi(2M_{n-1})=M_{n-1}+o(M_{n-1}).
\]
Letting $\zeta\downarrow0$ proves $|C_n|/M_{n-1}\to1$.
Taking successive quotients gives the third limit in \eqref{eq:weak-limits},
and dividing by $M_n/M_{n-1}\to\ph$ gives the fourth.

Finally, prefix completeness and \eqref{eq:prime-lower} imply
$Q_n\ge(1-\zeta)^2M_n$ eventually, for every $0<\zeta<1$;
indeed choose a prime between $(1-\zeta)^2M_n$ and $(1-\zeta)M_n$.
Together with $Q_n\le M_n$, this proves $Q_n/M_n\to1$.
\end{proof}

\appendix
\section{Why the old summability argument cannot simply be weakened}

The previous proof constructs cutoffs with relative losses summable along
exponentially growing scales, and deduces a \emph{positive} limiting
amplitude. Little-$o$ relative losses need not be summable. For example,
\[
 x_n=\frac{\ph^n}{n+3}
\]
is increasing and obeys the Fibonacci upper bound. For $n\ge1$, writing
$m=n+3$, direct calculation gives
\[
 \frac{x_n+x_{n-1}-x_{n+1}}{x_n}
   =\frac1{\ph(m-1)}+\frac\ph{m+1}=O(1/n)\longrightarrow0.
\]
Nevertheless $x_n/\ph^n=1/(n+3)\to0$. This example explains why
asymptotic Fibonacci saturation does not itself give a positive amplitude.
Our bounded-step proof obtains saturation first and then proves the ratio
limit directly, avoiding any such amplitude claim.

\section{Analytic and formalization boundary}

The two restricted hypotheses \eqref{eq:binary-weak} are of the same
fixed-interval type, with coefficients $(1,1)$ and $(2,1)$. Buffered
filling is derived from the first one, rather than imported as a
Conway-specific consequence of an almost-equal-primes theorem.

The preserved proof's circle-method appendix treats $(2,1)$ with stronger
logarithmic savings, using \cite[Proposition~24]{Tao} and
\cite[Lemma~1]{Kumchev}. The parallel $(1,1)$ calculation uses $S_I(\theta)$
instead of $S_I(2\theta)$ and even targets. Its major-arc coefficient is
$\mu(q)^2/\tot(q)^2$, and its real integral is
$\meas\{v\in I_X:N-v\in J_X\}$. For even $N$, the factor at $2$
is again $2$: now $c_2(N)=1$, and the coefficient for $q=2r$ equals
the odd-$r$ coefficient. All odd-prime local factors are the same as
in that appendix. Thus the standard analytic estimates there also give
the stronger fixed-interval $(1,1)$ result, and hence the weaker input
used here. This observation identifies the classical analytic source;
it is not a claim that either restricted theorem is already available
as a proved Lean declaration.

For a Lean implementation the revised target can be organized around
$\operatorname{Pref}(j,L)$ alone: neither $Q_n$ nor a first missing prime
is needed for the main ratio argument. The substantive new combinatorial
lemma is Lemma~\ref{lem:boost}; the new real-sequence argument is the
finite-step comparison in Proposition~\ref{prop:prefix-close}.
Natural-density-zero errors alone remain outside this reduction.
A proof using only those weaker errors would need an additional argument
showing that they cannot concentrate on the prime candidates.

\begin{thebibliography}{9}
\bibitem{CVZ}
M. Caragiu, P. A. Vicol, and M. Zaki,
\href{https://www.fq.math.ca/Papers1/55-4/CaragiuVicolZaki03162017.pdf}
{On Conway's subprime function, a covering of $\N$ and an unexpected
appearance of the golden ratio},
\emph{The Fibonacci Quarterly} \textbf{55} (2017), no.~4, 327--331.
\bibitem{Tao}
T. Tao,
\href{https://terrytao.wordpress.com/2015/03/30/254a-notes-8-the-hardy-littlewood-circle-method-and-vinogradovs-theorem/}
{254A, Notes 8: The Hardy--Littlewood circle method and Vinogradov's theorem},
30 March 2015, especially Proposition~24.
\bibitem{Kumchev}
A. V. Kumchev,
\href{https://tigerweb.towson.edu/akumchev/a23.pdf}
{On sums of primes from Beatty sequences},
\emph{Integers} \textbf{8} (2008), Article A08; Lemma~1.
\end{thebibliography}
\end{document}
